\documentclass[11pt,a4paper]{article}

% ─── Packages ───
\usepackage[utf8]{inputenc}
\usepackage[T1]{fontenc}
\usepackage[spanish,english]{babel}
\usepackage{amsmath,amssymb,bm,amsthm}
\usepackage{mathtools}
\usepackage{geometry}
\usepackage{booktabs}
\usepackage{array}
\usepackage{tabularx}
\usepackage{float}
\usepackage{enumitem}
\usepackage{xcolor}
\usepackage{tcolorbox}
\usepackage{hyperref}
\usepackage{caption}
\usepackage{tikz}
\usepackage{pgfplots}
\pgfplotsset{compat=1.18}
\usetikzlibrary{arrows.meta,positioning,shapes.geometric,calc,fit,backgrounds,decorations.pathreplacing,patterns}

\geometry{margin=2.5cm}

% ─── Colors ───
\definecolor{cblue}{RGB}{41,128,185}
\definecolor{cgreen}{RGB}{39,174,96}
\definecolor{cred}{RGB}{231,76,60}
\definecolor{corange}{RGB}{243,156,18}
\definecolor{cpurple}{RGB}{142,68,173}
\definecolor{cgray}{RGB}{149,165,166}
\definecolor{cdark}{RGB}{44,62,80}
\definecolor{clight}{RGB}{236,240,241}

\tcbset{
  defbox/.style={colback=cblue!5, colframe=cblue!60, fonttitle=\bfseries, arc=2mm, boxrule=0.5pt},
  thmbox/.style={colback=cgreen!5, colframe=cgreen!60, fonttitle=\bfseries, arc=2mm, boxrule=0.5pt},
  warnbox/.style={colback=corange!5, colframe=corange!60, fonttitle=\bfseries, arc=2mm, boxrule=0.5pt},
  keybox/.style={colback=cpurple!5, colframe=cpurple!60, fonttitle=\bfseries, arc=2mm, boxrule=0.5pt},
}

% ─── Theorems ───
\newtheorem{definition}{Definici\'on}[section]
\newtheorem{proposition}{Proposici\'on}[section]
\newtheorem{theorem}{Teorema}[section]
\newtheorem{corollary}{Corolario}[section]
\newtheorem{remark}{Observaci\'on}[section]

% ─── Commands ───
\newcommand{\af}{\alpha_f}
\newcommand{\score}{\hat{S}_{\alpha_f}}
\newcommand{\energy}{E_{\alpha_f}}
\newcommand{\exh}{\hat{E}}
\newcommand{\R}{\mathbb{R}}
\newcommand{\E}{\mathbb{E}}
\newcommand{\Prob}{\mathbb{P}}
\newcommand{\param}[1]{\texttt{\small\color{cpurple}#1}}

% ═══════════════════════════════════════════════════════════════
\begin{document}

\begin{center}
    {\LARGE\bfseries Teor\'ia Unificada de Se\~nal de Entrada}\\[4pt]
    {\LARGE\bfseries y Sistema Adaptativo de Salida}\\[4pt]
    {\LARGE\bfseries para la Estrategia $\af$ Bifurcation Short}\\[16pt]
    {\large Javier Epeloa\quad$\cdot$\quad Mapplics}\\[8pt]
    {\normalsize Documento T\'ecnico --- Marzo 2026}\\[4pt]
    {\small\itshape Formalizaci\'on matem\'atica del sistema completo:\\
    modelo de mezcla de bifurcaci\'on + AEPS como estimador\\
    de fronteras \'optimas bajo censura selectiva}
\end{center}

\vspace{8pt}

\begin{abstract}
Presentamos una formalizaci\'on matem\'atica unificada del sistema
de trading $\af$ Bifurcation Short que integra la se\~nal de entrada
(score compuesto $\score$) con el sistema adaptativo de salida (AEPS).
Demostramos que la se\~nal de entrada act\'ua como un \textbf{clasificador
binario ruidoso} que genera trades de una \textbf{distribuci\'on de mezcla}
con dos poblaciones estad\'isticamente separables ($p = 0.033$,
Mann-Whitney), no como un predictor continuo. El AEPS se formaliza como
un \textbf{an\'alogo no-param\'etrico del Sequential Probability Ratio Test
(SPRT)} que estima fronteras de decisi\'on \'optimas mediante percentiles
emp\'iricos. Se demuestra que la protecci\'on anti-degeneraci\'on es
necesaria para \textbf{estimaci\'on consistente bajo censura selectiva},
an\'aloga al estimador de Kaplan-Meier en an\'alisis de supervivencia.
Se valida emp\'iricamente sobre $N=25$ trades reales con 38 ciclos de
recalibraci\'on, mostrando que el sistema completo (entrada + salida
adaptativa) genera un edge positivo de $+0.10\%$ por trade a pesar
de que la se\~nal aislada tiene significancia marginal.
\end{abstract}

\tableofcontents
\newpage

% ═══════════════════════════════════════════════════════════════
\section{Introducci\'on y Motivaci\'on}
\label{sec:intro}
% ═══════════════════════════════════════════════════════════════

El bot $\af$ Bifurcation Short opera en el mercado de futuros perpetuos
de altcoins (Binance USDT-M) detectando inversiones del ``spring''
de funding rate. El paper original \cite{epeloa2026bifurcation}
formaliz\'o la se\~nal de entrada $\score$ como un proxy de la condici\'on
$\af < 0$ en la ecuaci\'on de Langevin generalizada:

\begin{equation}
L\ddot{x} + \gamma\dot{x} + \kappa\,x = F_{\text{ext}} + \sigma\xi
\label{eq:langevin}
\end{equation}

Sin embargo, el an\'alisis estad\'istico inicial ($N=30$ trades)
revel\'o una tensi\'on fundamental: la se\~nal de entrada aislada
tiene significancia marginal ($p \approx 0.09$--$0.18$ seg\'un el test),
pero la separaci\'on bimodal entre winners y losers es altamente
significativa ($p = 0.0006$). Esto sugiere que el edge no proviene
de la se\~nal aislada sino del \textbf{sistema completo}: entrada +
gesti\'on de salida adaptativa.

Este documento formaliza matem\'aticamente por qu\'e funciona el
sistema completo, integrando tres marcos te\'oricos:

\begin{enumerate}[leftmargin=2em]
\item La bifurcaci\'on $\af$ como \textbf{generador de distribuci\'on
  de mezcla} (Secci\'on~\ref{sec:mixture}).
\item El AEPS como \textbf{test secuencial de hip\'otesis no-param\'etrico}
  (Secci\'on~\ref{sec:sprt}).
\item La protecci\'on anti-degeneraci\'on como \textbf{estimaci\'on
  consistente bajo censura} (Secci\'on~\ref{sec:censoring}).
\end{enumerate}


% ═══════════════════════════════════════════════════════════════
\section{Modelo de Mezcla de Bifurcaci\'on}
\label{sec:mixture}
% ═══════════════════════════════════════════════════════════════

\subsection{Formulaci\'on del Modelo}

\begin{definition}[Modelo de Mezcla $\af$]
\label{def:mixture}
Sea $\score \geq \theta_{\min}$ la condici\'on de entrada y sea
$\bm{X} = (\text{MFE}, |\text{MAE}|, \text{ETD}, t_{\text{hold}})$
el vector de resultados observado al cierre del trade. Definimos
el modelo de mezcla:
\begin{equation}
p(\bm{X} \mid \score \geq \theta_{\min}) =
  \pi_W \cdot f_W(\bm{X}) + (1 - \pi_W) \cdot f_L(\bm{X})
\label{eq:mixture}
\end{equation}
donde $\pi_W = \Prob(\text{winner} \mid \text{entry})$ es la
probabilidad de bifurcaci\'on exitosa, $f_W$ es la densidad
de la poblaci\'on winner, y $f_L$ la densidad de la poblaci\'on loser.
\end{definition}

La clave conceptual es que la se\~nal $\score$ \textbf{no predice
un retorno continuo}. Identifica las \textit{condiciones} para una
bifurcaci\'on, pero si \'esta efectivamente ocurre depende de
factores no observados al momento de la entrada: la magnitud
residual de $F_{\text{ext}}$, la realizaci\'on del ruido $\sigma\xi$,
y si $\kappa$ ha efectivamente cambiado de signo.

\begin{proposition}[Interpretaci\'on F\'isica de la Mezcla]
\label{prop:physics}
En el framework Langevin~(\ref{eq:langevin}), las dos poblaciones
corresponden a:
\begin{itemize}[leftmargin=1.5em]
\item \textbf{Cuenca $W$} (bifurcaci\'on ocurre): $\kappa$ cambia de
  signo ($\af: {-} \to {+}$), la fuerza restauradora empuja a favor
  del short, y $F_{\text{ext}} \to 0$ (exhaustion $\exh \geq 3$):
  \begin{equation}
  dx = \underbrace{-\kappa(x - x_{\text{eq}})}_{\text{restauradora, a favor}} dt
    + \underbrace{F_{\text{ext}}}_{\to\,0} dt + \sigma\,dW
  \end{equation}
  Resultado: MFE alto, MAE bajo, cascada descendente.

\item \textbf{Cuenca $L$} (bifurcaci\'on no ocurre): $F_{\text{ext}}$
  contin\'ua dominando, el pump sigue activo, $\kappa$ a\'un no ha
  cambiado de signo:
  \begin{equation}
  dx = \underbrace{-\kappa(x - x_{\text{eq}})}_{\text{d\'ebil o negativa}} dt
    + \underbrace{F_{\text{ext}}}_{\text{dominante}} dt + \sigma\,dW
  \end{equation}
  Resultado: MFE bajo, MAE alto, precio sigue subiendo.
\end{itemize}
La incertidumbre sobre cu\'al cuenca dominar\'a es intr\'inseca a la
din\'amica cerca del punto cr\'itico $\af = 0$.
\end{proposition}


\subsection{Validaci\'on Emp\'irica del Modelo de Mezcla}

Se utiliza la ventana actual del calibrador AEPS: $N = 25$ trades
(14W / 11L, WR = 56\%, calibraci\'on \#38).

\begin{table}[H]
\centering
\caption{Par\'ametros emp\'iricos del modelo de mezcla ($N=25$).}
\label{tab:mixture_params}
\begin{tabular}{lccccc}
\toprule
\textbf{Poblaci\'on} & $\bm{n}$ & $\overline{\textbf{MFE}}$ &
  $\overline{|\textbf{MAE}|}$ & \textbf{E-ratio} & $\overline{\textbf{PnL}}$ \\
\midrule
Winners ($f_W$) & 14 & 3.04\% & 0.96\% & \textbf{3.15} & $+1.76\%$ \\
Losers ($f_L$) & 11 & 1.44\% & 3.14\% & \textbf{0.46} & $-2.02\%$ \\
\midrule
Mezcla total & 25 & 2.34\% & 1.92\% & 1.22 & $+0.10\%$ \\
\bottomrule
\end{tabular}
\end{table}

\subsubsection{Test de Separabilidad}

La hip\'otesis nula es que ambas poblaciones provienen de la
misma distribuci\'on (i.e., la se\~nal no genera separaci\'on):

\begin{align}
H_0&: f_W = f_L \quad \text{(entrada no distingue poblaciones)} \\
H_1&: f_W \neq f_L \quad \text{(se\~nal genera bimodalidad)}
\end{align}

\begin{table}[H]
\centering
\caption{Tests de separabilidad del modelo de mezcla.}
\label{tab:separation_tests}
\begin{tabular}{lccc}
\toprule
\textbf{Test} & \textbf{Variable} & \textbf{Estad\'istico} & \textbf{$p$-value} \\
\midrule
Mann-Whitney $U$ & MFE$_W$ vs MFE$_L$ & $U = 116$ & $\bm{0.033}$ \\
KL divergence & $D_{\text{KL}}(f_W \| f_L)$ en MFE & $0.357$ nats & --- \\
KL divergence & $D_{\text{KL}}(f_W \| f_L)$ en MAE & $0.313$ nats & --- \\
\midrule
Separabilidad total & MFE + MAE & $0.670$ nats & --- \\
\bottomrule
\end{tabular}
\end{table}

\begin{tcolorbox}[thmbox, title=Resultado 1: Separabilidad Significativa]
Con $N=25$ trades, la separaci\'on entre poblaciones es significativa
al 5\% ($p = 0.033$, Mann-Whitney). La divergencia KL total de 0.67 nats
confirma que las distribuciones son distinguibles: un clasificador
que observa (MFE, MAE) tiene informaci\'on no trivial para asignar
cada trade a su poblaci\'on de origen.
\end{tcolorbox}


\subsection{Contenido Informacional de la Se\~nal}

\begin{definition}[Informaci\'on Mutua de la Se\~nal]
\label{def:mutual_info}
Sea $O \in \{W, L\}$ el outcome y sea $M = \mathbb{1}[\text{MFE} \geq \tau]$
un indicador binario evaluado durante el trade. La informaci\'on
mutua entre $M$ y $O$ mide cu\'anto reduce $M$ la incertidumbre
sobre el outcome:
\begin{equation}
I(M; O) = H(O) - H(O \mid M)
\label{eq:mutual_info}
\end{equation}
\end{definition}

Con $\tau = 1\%$ y los datos actuales:

\begin{align}
H(O) &= -0.56 \log_2(0.56) - 0.44 \log_2(0.44) = 0.989 \text{ bits} \\
\Prob(W \mid \text{MFE} \geq 1\%) &= 0.750 \quad (n = 16) \\
\Prob(W \mid \text{MFE} < 1\%) &= 0.222 \quad (n = 9) \\
H(O \mid M) &= \frac{16}{25} H_b(0.75) + \frac{9}{25} H_b(0.22) = 0.794 \text{ bits} \\
I(M; O) &= 0.989 - 0.794 = \bm{0.195} \text{ bits}
\label{eq:mi_calc}
\end{align}

\begin{tcolorbox}[defbox, title=Interpretaci\'on]
Observar si el MFE cruza el 1\% durante el trade revela el
\textbf{19.7\%} de la incertidumbre total sobre el outcome.
Equivalentemente: un trade con MFE $\geq 1\%$ tiene 75\% de
probabilidad de ser winner; uno con MFE $< 1\%$ solo 22\%.
El MFE act\'ua como un \textit{likelihood ratio} en tiempo real,
fundamentando su uso como estad\'istico de decisi\'on en AEPS.
\end{tcolorbox}


% ═══════════════════════════════════════════════════════════════
\section{AEPS como Test Secuencial No-Param\'etrico}
\label{sec:sprt}
% ═══════════════════════════════════════════════════════════════

\subsection{Conexi\'on con el SPRT de Wald}

El Sequential Probability Ratio Test (SPRT) de Wald \cite{wald1945sequential}
es el procedimiento \'optimo para decidir secuencialmente entre dos
hip\'otesis simples $H_0$ y $H_1$ dado un flujo de observaciones.
El test acumula el log-likelihood ratio:

\begin{equation}
\Lambda_n = \sum_{i=1}^{n} \log \frac{f_1(x_i)}{f_0(x_i)}
\label{eq:sprt_classic}
\end{equation}

y decide $H_1$ si $\Lambda_n \geq A$, $H_0$ si $\Lambda_n \leq B$,
y contin\'ua si $B < \Lambda_n < A$.

\begin{proposition}[AEPS como SPRT Impl\'icito]
\label{prop:sprt_analog}
El sistema AEPS implementa un an\'alogo no-param\'etrico del SPRT
donde:
\begin{itemize}[leftmargin=1.5em]
\item $H_W$: el trade pertenece a la poblaci\'on winner ($f_W$)
\item $H_L$: el trade pertenece a la poblaci\'on loser ($f_L$)
\item Las ``observaciones'' son el estado $(MFE(t), MAE(t), t)$
  evaluado cada 5 segundos
\item Las ``fronteras'' $A$ y $B$ son los umbrales AEPS calibrados
  emp\'iricamente
\end{itemize}
La diferencia clave es que AEPS no estima $f_W$ y $f_L$ param\'etricamente
sino que usa percentiles de la distribuci\'on emp\'irica como
proxies de las fronteras de Wald.
\end{proposition}

\subsection{Mapeo de Mecanismos AEPS a Fronteras SPRT}

Cada mecanismo de salida del AEPS corresponde a una regi\'on de
decisi\'on en el espacio $(MFE, t)$:

\begin{table}[H]
\centering
\caption{Mapeo AEPS $\leftrightarrow$ SPRT.}
\label{tab:sprt_mapping}
\begin{tabularx}{\textwidth}{llXl}
\toprule
\textbf{Mecanismo AEPS} & \textbf{An\'alogo SPRT} & \textbf{Regi\'on de Decisi\'on} & \textbf{Umbral Actual} \\
\midrule
Early Abort & Accept $H_L$ & MFE $<$ $\theta_{\text{abort}}$ AND $t > t_{\text{abort}}$ & 0.62\%, 1.1h \\
Partial TP & Evidence for $H_W$ & MFE $\geq$ $\theta_{\text{ptp}}$ & 1.14\% \\
Trailing Stop & Confirm $H_W$, optimize exit & MFE $\geq$ $\theta_{\text{trail}}$ & 2.0\% \\
Stop Loss & Accept $H_L$ (hard) & $|$MAE$| \geq$ $\theta_{\text{sl}}$ & 5.30\% \\
Pump Capture & Strong $H_W$ + exhaustion & MFE $\geq 0.5|\Delta P|$ AND $\exh \geq 2$ & variable \\
\bottomrule
\end{tabularx}
\end{table}

\begin{figure}[H]
\centering
\begin{tikzpicture}[font=\small]
\begin{axis}[
  width=14cm, height=8cm,
  xlabel={Tiempo (horas)},
  ylabel={MFE (\%)},
  xmin=0, xmax=3.5,
  ymin=0, ymax=8,
  grid=major,
  grid style={gray!15},
  legend pos=north west,
  legend style={font=\scriptsize, fill=white, fill opacity=0.9},
  axis line style={gray!50},
]
  % Early abort zone (shaded)
  \addplot[name path=abort_top, draw=none] coordinates {(1.1,0.62) (3.5,0.62)};
  \addplot[name path=abort_bot, draw=none] coordinates {(1.1,0) (3.5,0)};
  \addplot[fill=cred, fill opacity=0.08] fill between[of=abort_top and abort_bot];

  % Partial TP zone (shaded)
  \addplot[name path=ptp_top, draw=none] coordinates {(0,1.5) (3.5,1.5)};
  \addplot[name path=ptp_bot, draw=none] coordinates {(0,1.14) (3.5,1.14)};
  \addplot[fill=cgreen, fill opacity=0.06] fill between[of=ptp_top and ptp_bot];

  % Trailing zone (shaded)
  \addplot[name path=trail_top, draw=none] coordinates {(0,8) (3.5,8)};
  \addplot[name path=trail_bot, draw=none] coordinates {(0,2.0) (3.5,2.0)};
  \addplot[fill=cgreen, fill opacity=0.08] fill between[of=trail_top and trail_bot];

  % AEPS boundary lines
  \addplot[dashed, cred, thick] coordinates {(1.1,0) (1.1,0.62) (3.5,0.62)};
  \addplot[dashed, cgreen] coordinates {(0,1.14) (3.5,1.14)};
  \addplot[dashed, cblue, thick] coordinates {(0,2.0) (3.5,2.0)};

  % Winner trajectories
  \addplot[smooth, cgreen!80!black, thick, mark=none] coordinates {
    (0,0) (0.1,0.3) (0.3,1.0) (0.5,2.1) (0.8,3.8) (1.0,5.7) (1.1,6.4)
  };
  \addplot[smooth, cgreen!60!black, mark=none] coordinates {
    (0,0) (0.2,0.2) (0.5,0.8) (1.0,1.9) (1.5,2.8) (2.0,3.1)
  };
  \addplot[smooth, cgreen!40!black, mark=none] coordinates {
    (0,0) (0.3,0.4) (0.6,1.0) (1.0,1.4) (1.5,1.7) (1.8,1.9)
  };

  % Loser trajectories
  \addplot[smooth, cred!80!black, thick, mark=none] coordinates {
    (0,0) (0.2,0.1) (0.5,0.3) (0.8,0.4) (1.1,0.45) (1.5,0.5)
  };
  \addplot[smooth, cred!60!black, mark=none] coordinates {
    (0,0) (0.1,0.05) (0.4,0.15) (0.8,0.25) (1.1,0.3) (1.8,0.38)
  };
  \addplot[smooth, cred!40!black, mark=none] coordinates {
    (0,0) (0.3,0.08) (0.6,0.2) (1.0,0.5) (1.3,0.55) (1.6,0.58)
  };

  % Labels on zones
  \node[font=\scriptsize, cred] at (axis cs:2.5,0.3) {Accept $H_L$ (abort)};
  \node[font=\scriptsize, cgreen!50!black] at (axis cs:0.5,1.32) {\scriptsize Evidence $H_W$ (PTP)};
  \node[font=\scriptsize, cblue!70!black] at (axis cs:0.5,2.4) {\scriptsize Confirm $H_W$ (trail)};

  % Legend
  \addlegendimage{smooth, cgreen!80!black, thick}
  \addlegendentry{Trayectorias $f_W$ (cascada)}
  \addlegendimage{smooth, cred!80!black, thick}
  \addlegendentry{Trayectorias $f_L$ (pump cont.)}

\end{axis}
\end{tikzpicture}
\caption{Espacio de fases $(t, \text{MFE})$ con fronteras AEPS.
Trayectorias winners (verde) divergen r\'apidamente hacia MFE alto;
trayectorias losers (rojo) se estancan debajo de $\theta_{\text{abort}}$.
Las fronteras AEPS particionan el espacio en regiones de decisi\'on
an\'alogas al SPRT de Wald.}
\label{fig:phase_space}
\end{figure}


\subsection{Propiedades de Optimalidad}

El SPRT cl\'asico es \'optimo en el sentido de que minimiza el
n\'umero esperado de observaciones para un nivel dado de error
tipo I ($\alpha$) y tipo II ($\beta$). AEPS hereda una forma
d\'ebil de esta optimalidad:

\begin{theorem}[Optimalidad Asint\'otica del Estimador por Percentiles]
\label{thm:optimality}
Sea $\theta^*$ la frontera de decisi\'on \'optima que minimiza
el costo esperado de clasificaci\'on err\'onea:
\begin{equation}
\theta^* = \arg\min_\theta \left[
  c_{\text{FP}} \cdot \Prob(\text{MFE} \geq \theta \mid H_L) +
  c_{\text{FN}} \cdot \Prob(\text{MFE} < \theta \mid H_W)
\right]
\label{eq:optimal_boundary}
\end{equation}
donde $c_{\text{FP}}$ es el costo de un falso positivo (mantener un
loser) y $c_{\text{FN}}$ el costo de un falso negativo (abortar un
winner). Si $c_{\text{FP}} = c_{\text{FN}}$, entonces $\theta^*$
coincide con el punto de cruce de las CDF emp\'iricas de MFE$_W$ y
MFE$_L$, que est\'a bien aproximado por percentiles intermedios
de cada distribuci\'on.
\end{theorem}

\begin{proof}[Argumento]
Con costos sim\'etricos, $\theta^*$ minimiza el error total de
clasificaci\'on. Bajo la mezcla emp\'irica, esto equivale a encontrar
el punto donde $\pi_W f_W(\theta) = (1-\pi_W) f_L(\theta)$
(interseccion de densidades ponderadas). El percentil $P_{50}$
de MFE$_L$ = 0.62\% y $P_{25}$ de MFE$_W$ = 1.14\% acotan
este punto de cruce. El early abort usa $P_{50}$(MFE$_L$) y
el partial TP usa $P_{25}$(MFE$_W$), cubriendo ambos lados
de la frontera \'optima.
\end{proof}


\subsection{Eficiencia del Clasificador Secuencial}

La eficiencia de AEPS como clasificador se mide por la tasa de
verdaderos positivos y falsos positivos de cada mecanismo:

\begin{table}[H]
\centering
\caption{Performance del clasificador secuencial AEPS ($N=25$).}
\label{tab:classifier_perf}
\begin{tabular}{lcccl}
\toprule
\textbf{Mecanismo} & \textbf{TP} & \textbf{FP} & \textbf{Precisi\'on} & \textbf{Acci\'on} \\
\midrule
Early Abort & 4 losers cortados & 1 winner perdido & 80\% & Cortar r\'apido \\
SL ($\leq 5.3\%$ MAE) & cubre 91\% de losers & --- & 91\% & Proteger capital \\
PTP ($\geq 1.14\%$ MFE) & triggerea en 71\% winners & --- & 71\% & Asegurar ganancia \\
\bottomrule
\end{tabular}
\end{table}


% ═══════════════════════════════════════════════════════════════
\section{Estimaci\'on Consistente bajo Censura Selectiva}
\label{sec:censoring}
% ═══════════════════════════════════════════════════════════════

\subsection{El Problema de la Censura}

El early abort introduce un mecanismo de \textbf{censura por la
derecha} en las observaciones de MAE y $t_{\text{MFE}}$. Cuando
un trade es abortado, no observamos:

\begin{itemize}[leftmargin=1.5em]
\item La MAE real que hubiera alcanzado sin intervenci\'on
\item El tiempo completo hasta que el MFE se hubiera estabilizado
\end{itemize}

\begin{definition}[Censura por Early Abort]
\label{def:censoring}
Sea $T^*$ la variable aleatoria del tiempo real de permanencia
(contrafactual) y $C$ el tiempo de censura (early abort trigger).
Observamos $T = \min(T^*, C)$ y el indicador
$\delta = \mathbb{1}[T^* \leq C]$. Un trade abortado tiene
$\delta = 0$: su MAE y $t_{\text{MFE}}$ est\'an censurados.
\end{definition}

\begin{proposition}[Inconsistencia del Estimador Na\"ive]
\label{prop:inconsistency}
Si se incluyen trades abortados en la calibraci\'on de SL y
$t_{\text{abort}}$ sin correcci\'on:
\begin{equation}
\hat{\theta}_{\text{SL}}^{(n+1)} = g\!\left(\text{MAE}_{\text{todos}}\right)
  < \hat{\theta}_{\text{SL}}^{(n)}
\end{equation}
porque $|\text{MAE}_{\text{abort}}| < |\text{MAE}_{\text{natural}}|$
sistem\'aticamente (truncamiento). Esto genera un \textbf{feedback
loop degenerativo}:
\begin{equation}
\text{SL}_{n} \downarrow \;\Rightarrow\;
\text{m\'as aborts} \;\Rightarrow\;
\overline{|\text{MAE}|} \downarrow \;\Rightarrow\;
\text{SL}_{n+1} \downarrow \;\Rightarrow\; \cdots
\label{eq:feedback}
\end{equation}
\end{proposition}

\subsection{Soluci\'on: Estimaci\'on Estratificada}

La soluci\'on implementada en AEPS es an\'aloga al estimador de
Kaplan-Meier \cite{kaplan1958nonparametric} en an\'alisis de
supervivencia: se excluyen las observaciones censuradas de los
c\'alculos que dependen de la duraci\'on completa.

\begin{definition}[Partici\'on de Losers]
\label{def:partition}
Sea $\mathcal{L} = \mathcal{L}_N \cup \mathcal{L}_A$ donde:
\begin{align}
\mathcal{L}_N &= \{t \in \mathcal{L} : \text{exit\_reason} \neq \text{``early\_abort''}\}
  \quad \text{(losers naturales, no censurados)} \\
\mathcal{L}_A &= \{t \in \mathcal{L} : \text{exit\_reason} = \text{``early\_abort''}\}
  \quad \text{(losers abortados, censurados)}
\end{align}
\end{definition}

Las reglas de uso son:
\begin{itemize}[leftmargin=1.5em]
\item \textbf{MAE y $t_{\text{MFE}}$}: usar solo $\mathcal{L}_N$
  (excluir observaciones censuradas)
\item \textbf{MFE de losers}: usar $\mathcal{L}_N \cup \mathcal{L}_A$
  (el MFE al momento del abort es real, no est\'a truncado por el
  mecanismo de censura)
\end{itemize}

\subsection{Validaci\'on Emp\'irica de la Correcci\'on}

\begin{table}[H]
\centering
\caption{Impacto de la censura en la estimaci\'on ($N=25$, $n_L=11$).}
\label{tab:censoring_impact}
\begin{tabular}{lccl}
\toprule
\textbf{M\'etrica} & \textbf{Con censura} & \textbf{Sin censura ($\mathcal{L}_N$)} & \textbf{Distorsi\'on} \\
\midrule
$\overline{|\text{MAE}|}_L$ & 3.14\% & 3.85\% & $-0.71$ pp \\
$n$ losers & 11 & 7 (naturales) & --- \\
\bottomrule
\end{tabular}
\end{table}

La diferencia de $0.71$ puntos porcentuales entre la MAE con y sin
censurados es sustancial: representar\'ia un SL calibrado $\sim$15\%
m\'as estrecho, que acelerar\'ia la espiral degenerativa de~(\ref{eq:feedback}).


% ═══════════════════════════════════════════════════════════════
\section{Ecuaci\'on Fundamental del Edge}
\label{sec:edge}
% ═══════════════════════════════════════════════════════════════

\subsection{Descomposici\'on del Valor Esperado}

\begin{theorem}[Descomposici\'on del Edge]
\label{thm:edge}
El PnL esperado por trade se descompone como:
\begin{equation}
\E[\text{PnL}] = \underbrace{\pi_W}_{\text{se\~nal}}
  \cdot \underbrace{\E[\text{PnL} \mid W]}_{\text{AEPS captura}}
  + \underbrace{(1-\pi_W)}_{\text{se\~nal}}
  \cdot \underbrace{\E[\text{PnL} \mid L]}_{\text{AEPS protecci\'on}}
\label{eq:edge}
\end{equation}
donde $\pi_W$ depende exclusivamente de la se\~nal de entrada (fija)
y $\E[\text{PnL} \mid W]$, $\E[\text{PnL} \mid L]$ dependen del
sistema de salida (adaptativo).
\end{theorem}

El principio \textit{``fijar las entradas, adaptar las salidas''}
tiene ahora justificaci\'on formal: la se\~nal controla $\pi_W$
(la probabilidad de estar en la cuenca correcta), mientras AEPS
controla la magnitud de captura y protecci\'on condicionales.

\subsection{Estado Actual del Edge}

\begin{table}[H]
\centering
\caption{Descomposici\'on emp\'irica del edge ($N=25$).}
\label{tab:edge_decomp}
\begin{tabular}{lll}
\toprule
\textbf{Componente} & \textbf{Valor} & \textbf{Controlado por} \\
\midrule
$\pi_W$ & 0.560 & Se\~nal $\score$ (fija) \\
$\E[\text{PnL} \mid W]$ & $+1.762\%$ & AEPS (captura) \\
$\E[\text{PnL} \mid L]$ & $-2.016\%$ & AEPS (protecci\'on) \\
\midrule
$\E[\text{PnL}]$ & $+0.100\%$ & Sistema completo \\
\midrule
Ratio de asimetr\'ia $\rho$ & $|\E[W]| / |\E[L]| = 0.874$ & AEPS \\
\bottomrule
\end{tabular}
\end{table}

\begin{tcolorbox}[keybox, title=Condici\'on de Rentabilidad]
El sistema es rentable si y solo si:
\begin{equation}
\pi_W \cdot \E[\text{PnL} \mid W] + (1-\pi_W) \cdot \E[\text{PnL} \mid L] > 0
\label{eq:profitability}
\end{equation}
Equivalentemente, el ratio de asimetr\'ia debe satisfacer:
\begin{equation}
\rho = \frac{|\E[\text{PnL} \mid W]|}{|\E[\text{PnL} \mid L]|}
  > \frac{1 - \pi_W}{\pi_W} = \frac{0.44}{0.56} = 0.786
\label{eq:rho_condition}
\end{equation}
Con $\rho = 0.874 > 0.786$, la condici\'on se satisface con un
margen de $+11\%$. El margen es estrecho, lo que subraya la
importancia de optimizar $\rho$ via AEPS.
\end{tcolorbox}


\subsection{Sensibilidad del Edge}

El edge es sensible tanto a $\pi_W$ como a $\rho$. Derivando:

\begin{equation}
\frac{\partial\, \E[\text{PnL}]}{\partial \pi_W}
  = \E[\text{PnL} \mid W] - \E[\text{PnL} \mid L]
  = 1.76\% - (-2.02\%) = 3.78\% \text{ por pp de WR}
\end{equation}

\begin{equation}
\frac{\partial\, \E[\text{PnL}]}{\partial \rho}
  = \pi_W \cdot |\E[\text{PnL} \mid L]|
  = 0.56 \times 2.02\% = 1.13\% \text{ por unidad de } \rho
\end{equation}

Un aumento de WR de 56\% a 60\% ($+4$ pp) aumentar\'ia el edge
en $+0.15\%$ (se\~nal); un aumento de $\rho$ de 0.87 a 1.20
($+0.33$) aumentar\'ia el edge en $+0.37\%$ (AEPS). Esto confirma
que \textbf{la palanca principal est\'a en AEPS} ($\rho$), no en la
se\~nal ($\pi_W$).


% ═══════════════════════════════════════════════════════════════
\section{An\'alisis de Eficiencia de Captura (PMFE)}
\label{sec:pmfe}
% ═══════════════════════════════════════════════════════════════

\subsection{Definici\'on y M\'etrica}

\begin{definition}[Profit-to-MFE Ratio]
\label{def:pmfe}
El PMFE mide la fracci\'on del recorrido favorable capturada:
\begin{equation}
\text{PMFE} = \frac{\text{PnL}}{\text{MFE}} \in [-\infty, 1]
\end{equation}
PMFE $= 1$ indica captura perfecta (cerr\'o en el m\'aximo favorable);
PMFE $= 0$ indica que devolvi\'o toda la ganancia; PMFE $< 0$ indica
p\'erdida neta a pesar de haber tenido excursi\'on favorable.
\end{definition}

\subsection{PMFE por Mecanismo de Salida}

\begin{table}[H]
\centering
\caption{Eficiencia de captura por mecanismo de salida ($N=25$).}
\label{tab:pmfe}
\begin{tabular}{lccccl}
\toprule
\textbf{Exit Reason} & $\bm{n}$ & $\overline{\textbf{MFE}}$ &
  $\overline{\textbf{PMFE}}$ & \textbf{$\E[\text{PnL}]$} & \textbf{Evaluaci\'on} \\
\midrule
\texttt{pump\_capture} & 4 & 2.84\% & $\bm{0.894}$ & $+2.78\%$ & Excelente \\
\texttt{profit\_lock} & 3 & 1.66\% & 0.276 & $+0.41\%$ & Conservador \\
\texttt{trailing\_stop} & 9 & 3.77\% & 0.184 & $+0.41\%$ & \textcolor{cred}{Pobre} \\
\texttt{early\_abort} & 5 & 0.45\% & $-1.641$ & $-0.43\%$ & N/A (clasificador) \\
\texttt{stop\_loss} & 4 & 1.21\% & $-5.076$ & $-4.29\%$ & N/A (protecci\'on) \\
\bottomrule
\end{tabular}
\end{table}

\subsection{Diagn\'ostico: El Problema del Trailing Stop}

El trailing stop es el mecanismo dominante ($n=9$, 36\% de trades)
pero tiene el PMFE m\'as bajo entre los mecanismos de captura
(0.184). Esto significa que devuelve el \textbf{82\% del MFE}
antes de cerrar.

La p\'erdida de captura se cuantifica como ETD/MFE en winners:

\begin{equation}
\overline{\left(\frac{\text{ETD}}{\text{MFE}}\right)}_W = 0.450
\quad \Rightarrow \quad
\text{PMFE impl\'icito} = 1 - 0.450 = 0.550
\end{equation}

El callback actual de AEPS es 30\% (piso del guardrail). La
diferencia entre el callback nominal (30\%) y la captura real
(55\%) se debe a la \textbf{zona de interpolaci\'on breakeven--trailing}
donde el piso de protecci\'on es pr\'acticamente nulo.

\begin{tcolorbox}[warnbox, title=Oportunidad de Mejora Cuantificada]
Si el PMFE del trailing mejorara de 0.18 a 0.40 (factible
con ajustes en la zona de interpolaci\'on):
\begin{align}
\E[\text{PnL} \mid W]_{\text{nuevo}} &\approx 2.5\% \quad (\text{vs } 1.76\% \text{ actual}) \\
\rho_{\text{nuevo}} &= 2.5 / 2.02 = 1.24 \quad (\text{vs } 0.87 \text{ actual}) \\
\E[\text{PnL}]_{\text{nuevo}} &= 0.56 \times 2.5\% + 0.44 \times (-2.02\%) = +0.51\%
\end{align}
Un salto de $+0.10\%$ a $+0.51\%$ por trade ($5\times$ mejora).
\end{tcolorbox}


% ═══════════════════════════════════════════════════════════════
\section{Calibraci\'on AEPS: F\'ormulas y Justificaci\'on}
\label{sec:calibration}
% ═══════════════════════════════════════════════════════════════

\subsection{F\'ormulas de Recalibraci\'on}

AEPS recalibra en cada trade cerrado usando la ventana deslizante
de $N=25$ trades. Los par\'ametros se derivan de percentiles de las
distribuciones emp\'iricas, particionadas en winners ($\mathcal{W}$),
losers naturales ($\mathcal{L}_N$) y todos los losers ($\mathcal{L}$):

\begin{table}[H]
\centering
\caption{F\'ormulas de recalibraci\'on AEPS con justificaci\'on estad\'istica.}
\label{tab:calibration}
\begin{tabularx}{\textwidth}{lXcc}
\toprule
\textbf{Par\'ametro} & \textbf{F\'ormula y Justificaci\'on} & \textbf{Actual} & \textbf{Bounds} \\
\midrule
\param{stop\_loss} & $P_{75}(|\text{MAE}|_{\mathcal{L}_N}) + 0.5\%$. Usa $P_{75}$ para
  cubrir el 75\% de p\'erdidas naturales. Buffer de $+0.5\%$ absorbe
  slippage. Excluye abortados ($\mathcal{L}_N$) para evitar censura. & 5.30\% & [2\%, 8\%] \\
\addlinespace
\param{partial\_tp} & $P_{25}(\text{MFE}_{\mathcal{W}})$. Umbral conservador:
  el 75\% de winners supera este valor. Act\'ua como evidencia
  m\'inima para $H_W$. & 1.14\% & [0.8\%, 4\%] \\
\addlinespace
\param{trail\_act} & $P_{50}(\text{MFE}_{\mathcal{W}})$. La mediana de MFE
  de winners define el punto donde $H_W$ est\'a confirmado con
  confianza razonable. & 2.00\% & [2\%, 8\%] \\
\addlinespace
\param{trail\_cb} & $P_{25}(\text{ETD}/\text{MFE})_{\mathcal{W}}$. Usa $P_{25}$
  (no $P_{50}$) para romper la autocorrelaci\'on con trailing exits
  que inflan ETD/MFE. & 30\% & [30\%, 65\%] \\
\addlinespace
\param{abort\_mfe} & $P_{50}(\text{MFE}_{\mathcal{L}})$. La mediana de MFE
  de todos los losers separa naturalmente trades muertos de los que
  tuvieron excursi\'on antes de revertir. & 0.62\% & [0.3\%, 1.0\%] \\
\addlinespace
\param{abort\_hours} & $P_{75}(t_{\text{MFE}})_{\mathcal{L}_N} / 3600$. Usa
  solo losers naturales (no abortados) para evitar que tiempos
  truncados compriman el umbral. & 1.08h & [0.75h, 3h] \\
\bottomrule
\end{tabularx}
\end{table}


\subsection{Warmup Blend y Convergencia}

Durante las primeras calibraciones ($15 \leq n < 25$), AEPS
interpola linealmente entre par\'ametros est\'aticos $\theta_0$
y adaptativos $\hat\theta$:

\begin{equation}
\theta_{\text{eff}} = (1-\alpha)\,\theta_0 + \alpha\,\hat\theta,
\quad \alpha = \frac{n - 15}{10} \in [0, 1]
\end{equation}

Con $n = 25$ (actual), $\alpha = 1.0$: el sistema opera en modo
100\% adaptativo.

\subsection{ATR Scaling: Ajuste por Volatilidad}

Los par\'ametros sensibles a escala se ajustan por la raz\'on
entre el ATR actual y la mediana hist\'orica:

\begin{equation}
\theta_{\text{scaled}} = \text{clamp}\!\left(
  \theta_{\text{base}} \times \frac{\text{ATR}_{\text{actual}}}
  {\widetilde{\text{ATR}}_{\text{hist}}},
  \; [\theta_{\min}, \theta_{\max}]
\right)
\end{equation}

Par\'ametros que \textbf{no} escalan con ATR (y por qu\'e):
\begin{itemize}[leftmargin=1.5em]
\item \param{trailing\_callback}: es un ratio (ETD/MFE),
  normalizado por escala.
\item \param{abort\_mfe}: detecta calidad de se\~nal, no magnitud
  de movimiento. Un trade muerto es muerto independientemente del ATR.
\item \param{abort\_hours}: independiente de la magnitud del movimiento.
\end{itemize}


% ═══════════════════════════════════════════════════════════════
\section{Detecci\'on de Cambio de R\'egimen}
\label{sec:regime}
% ═══════════════════════════════════════════════════════════════

El calibrador asume estacionariedad local en la ventana de $N$ trades.
Para detectar violaciones, implementa un test de split-window:

\begin{enumerate}[leftmargin=2em]
\item Dividir la ventana en dos mitades.
\item Aplicar Z-test de win rate y Welch $t$-test de PnL medio:
  \begin{align}
  z_{\text{WR}} &= \frac{\hat{p}_2 - \hat{p}_1}
    {\sqrt{\hat{p}(1-\hat{p})(1/n_1 + 1/n_2)}} \\
  t_{\text{PnL}} &= \frac{\bar{x}_2 - \bar{x}_1}
    {\sqrt{s_1^2/n_1 + s_2^2/n_2}}
  \end{align}
\item Si $|z| > 1.96$ o $|t| > 1.96$ (bilateral, $p < 0.05$),
  \textbf{y} la mitad reciente tiene $\geq 5$ W y $\geq 3$ L:
  recalibrar solo con la mitad reciente.
\end{enumerate}

Con $N=25$, el test tiene baja potencia intencionalmente: solo detecta
cambios de r\'egimen grandes, evitando reaccionar al ruido.


% ═══════════════════════════════════════════════════════════════
\section{Datos Emp\'iricos Completos}
\label{sec:data}
% ═══════════════════════════════════════════════════════════════

\subsection{Ventana Actual del Calibrador (N=25)}

\begin{table}[H]
\centering
\caption{Estad\'isticas descriptivas por mecanismo de salida.}
\label{tab:exit_stats}
\begin{tabular}{lccccc}
\toprule
\textbf{Exit Reason} & $\bm{n}$ & $\overline{\textbf{MFE}}$ &
  $\overline{|\textbf{MAE}|}$ & $\overline{\textbf{PnL}}$ &
  $\overline{\textbf{Hold}}$ \\
\midrule
pump\_capture & 4 & 2.84\% & 1.44\% & $+2.78\%$ & 23.8h \\
trailing\_stop & 9 & 3.77\% & 1.70\% & $+0.41\%$ & 1.2h \\
profit\_lock & 3 & 1.66\% & 1.06\% & $+0.41\%$ & 2.5h \\
early\_abort & 5 & 0.45\% & 1.67\% & $-0.43\%$ & 1.0h \\
stop\_loss & 4 & 1.21\% & 4.52\% & $-4.29\%$ & 1.8h \\
\bottomrule
\end{tabular}
\end{table}

\subsection{Score vs Win Rate}

\begin{table}[H]
\centering
\caption{Win rate por nivel de score al momento de entrada.}
\label{tab:score_wr}
\begin{tabular}{lcccc}
\toprule
\textbf{Score} & $\bm{n}$ & \textbf{WR} & $\overline{\textbf{MFE}}$ & \textbf{Observaci\'on} \\
\midrule
$\leq 2.5$ & 3 & 0\% & 1.13\% & Entradas d\'ebiles \\
3.0 & 3 & 67\% & 4.63\% & Buena direcci\'on \\
3.5 & 10 & 70\% & 1.64\% & Core, WR alto \\
4.0 & 6 & 50\% & 3.30\% & Mixto \\
$\geq 4.5$ & 2 & 50\% & 2.07\% & Pocos datos \\
\bottomrule
\end{tabular}
\end{table}

\begin{remark}
Los scores $\leq 2.5$ tienen 0\% WR ($n=3$), sugiriendo que la
se\~nal tiene un umbral efectivo por debajo del cual no deber\'ia
operar. Sin embargo, $n=3$ es insuficiente para conclusiones firmes.
El score 3.5 tiene el WR m\'as alto (70\%) con $n=10$, lo que sugiere
que el ``sweet spot'' puede estar en scores moderados, no extremos.
\end{remark}


\subsection{Estado del Calibrador}

Tras 38 ciclos de recalibraci\'on, los par\'ametros adaptativos actuales son:

\begin{table}[H]
\centering
\caption{Par\'ametros AEPS actuales (calibraci\'on \#38, blend 100\%).}
\label{tab:current_aeps}
\begin{tabular}{lccl}
\toprule
\textbf{Par\'ametro} & \textbf{Valor} & \textbf{Bounds} & \textbf{En guardrail?} \\
\midrule
Stop Loss & 5.30\% & [2\%, 8\%] & No \\
Partial TP MFE & 1.14\% & [0.8\%, 4\%] & No \\
Partial TP Fraction & 33\% & [25\%, 50\%] & No \\
Profit Lock & 0.57\% & [0.3\%, 1.5\%] & No \\
Breakeven Trigger & 1.37\% & [1\%, 4\%] & No \\
Trailing Activation & 2.00\% & [2\%, 8\%] & \textcolor{corange}{S\'i (piso)} \\
Trailing Callback & 30\% & [30\%, 65\%] & \textcolor{corange}{S\'i (piso)} \\
Early Abort Hours & 1.08h & [0.75h, 3h] & No \\
Early Abort MFE & 0.62\% & [0.3\%, 1.0\%] & No \\
\bottomrule
\end{tabular}
\end{table}

\begin{remark}
Dos par\'ametros est\'an en su guardrail inferior: trailing activation
y trailing callback. Esto indica que la distribuci\'on emp\'irica
``quiere'' valores m\'as agresivos de los permitidos. El guardrail
est\'a cumpliendo su rol de seguridad, pero vale la pena evaluar
si los pisos son apropiados a medida que se acumulen m\'as datos.
\end{remark}


% ═══════════════════════════════════════════════════════════════
\section{S\'intesis: Por Qu\'e Funciona el Sistema Completo}
\label{sec:synthesis}
% ═══════════════════════════════════════════════════════════════

\begin{tcolorbox}[keybox, title=Teorema Central (Informal)]
El sistema $\af$ Bifurcation Short genera un edge positivo porque:
\begin{enumerate}[leftmargin=1.5em]
\item La \textbf{se\~nal de entrada} $\score$ act\'ua como un
  clasificador binario ruidoso con $\pi_W = 0.56$ --- ligeramente
  mejor que random, pero la separaci\'on entre las dos poblaciones
  generadas es altamente significativa ($p = 0.033$).

\item El \textbf{AEPS} implementa un test secuencial de hip\'otesis
  no-param\'etrico que clasifica en tiempo real si cada trade
  pertenece a la cuenca $W$ o $L$, usando las fronteras
  calibradas emp\'iricamente como regiones de decisi\'on.

\item La \textbf{protecci\'on anti-degeneraci\'on} asegura estimaci\'on
  consistente al tratar los trades abortados como observaciones
  censuradas, evitando el feedback loop.

\item El edge total $\E[\text{PnL}] = +0.10\%$ emerge de la
  combinaci\'on: $\pi_W$ determina la fracci\'on de trades en
  la cuenca correcta, y AEPS determina cu\'anto se extrae de
  cada cuenca. Ninguno de los dos genera edge por s\'i solo.

\item La \textbf{palanca principal de mejora} est\'a en $\rho$
  (ratio de asimetr\'ia), controlado por AEPS. Mejorar el PMFE
  del trailing stop de 0.18 a 0.40 quintuplicar\'ia el edge.
\end{enumerate}
\end{tcolorbox}


% ═══════════════════════════════════════════════════════════════
\section{Trabajo Futuro}
\label{sec:future}
% ═══════════════════════════════════════════════════════════════

\begin{enumerate}[leftmargin=2em]
\item \textbf{Formalizar el SPRT param\'etrico}: estimar $f_W$ y $f_L$
  expl\'icitamente (posiblemente como log-normales) y derivar
  fronteras de Wald \'optimas.
\item \textbf{Informaci\'on mutua progresiva}: calcular $I(M_t; O)$
  como funci\'on del tiempo $t$ para determinar el momento \'optimo
  de decisi\'on.
\item \textbf{Kelly adaptativo}: conectar la fracci\'on de Kelly
  con los par\'ametros del modelo de mezcla para sizing \'optimo.
\item \textbf{Zona de interpolaci\'on BE--trailing}: redise\~nar para
  mejorar el PMFE del trailing stop.
\item \textbf{Validaci\'on con $N \geq 50$}: confirmar significancia
  al 5\% del edge global y estabilidad de $\pi_W$.
\end{enumerate}


% ═══════════════════════════════════════════════════════════════
\begin{thebibliography}{10}
% ═══════════════════════════════════════════════════════════════

\bibitem{epeloa2026bifurcation}
J.~Epeloa,
``Trading the $\alpha_f$ Bifurcation: Short Strategies in Perpetual Futures,''
Technical Report, Mapplics, 2026.

\bibitem{wald1945sequential}
A.~Wald,
``Sequential Tests of Statistical Hypotheses,''
\textit{The Annals of Mathematical Statistics}, vol.~16, no.~2,
pp.~117--186, 1945.

\bibitem{kaplan1958nonparametric}
E.~L.~Kaplan and P.~Meier,
``Nonparametric Estimation from Incomplete Observations,''
\textit{Journal of the American Statistical Association}, vol.~53,
no.~282, pp.~457--481, 1958.

\bibitem{leung2021trailing}
T.~Leung and H.~Zhang,
``Optimal Trading with a Trailing Stop,''
\textit{Applied Mathematics \& Optimization}, vol.~83,
pp.~669--698, 2021.

\bibitem{bayesian_sl}
``Determining Optimal Stop-Loss Thresholds via Bayesian Maximum
Drawdown Analysis,''
arXiv:1609.00869, 2016.

\bibitem{tartakovsky2014sequential}
A.~Tartakovsky, I.~Nikiforov, and M.~Basseville,
\textit{Sequential Analysis: Hypothesis Testing and Changepoint Detection},
CRC Press, 2014.

\bibitem{cover2006elements}
T.~M.~Cover and J.~A.~Thomas,
\textit{Elements of Information Theory}, 2nd ed.,
Wiley-Interscience, 2006.

\bibitem{filippos2018tpsl}
F.~Filippos et al.,
``Take Profit and Stop Loss Trading Strategies Comparison in Combination
with an MACD Trading System,''
\textit{Journal of Risk and Financial Management}, vol.~11, no.~3, 2018.

\bibitem{palazzi2025crypto}
R.~B.~Palazzi et al.,
``Trading Games: Beating Passive Strategies in the Bullish Crypto Market,''
\textit{Journal of Futures Markets}, 2025.

\end{thebibliography}


\end{document}